% Paper 3, §14 — Contradictions live and open frontiers
% Anchors: kb/index/contradictions.md § reasoning, plus the body
%          \begin{contradiction} blocks of §02–§13.
% Cross-paper: Paper-4 §faithfulness and Paper-5 §AIME-contamination
%              both backreference §11 and §14 of this paper.

\section{Contradictions live and open frontiers}
\label{sec:contradictions}

The previous thirteen sections ran a leg-by-leg decomposition of the
reasoning triangle and flagged tensions as they arose. This section
collapses those tensions into one place. The role is not to relitigate
each result but to catalogue \emph{which} questions the 2026-Q2
literature has not closed and \emph{what shape} a closure would have
to take. Six contradictions are load-bearing for the rest of the
paper's claims; each is sourced from the body
\texttt{contradiction} blocks of \cref{sec:process-reward-models,sec:tree-search,sec:rlvr-grpo-dapo,sec:r1-family,sec:search-vs-rl,sec:cot-faithfulness,sec:prm-survival,sec:frontier-2026}
and from \texttt{kb/index/contradictions.md}~\S~reasoning.

The structural reading: the three legs of the triangle are real and
co-designed in frontier 2026 reasoning models. Where the literature
disagrees is not on whether the legs exist but on \emph{which
combination} is optimal at which $(C_\text{train}, C_\text{RL},
C_\text{test})$ point and at which prompt-difficulty regime. None of
the six contradictions has a closed-form answer in 2026-Q2, and none
is on track to acquire one before the next model generation. The
implication for the paper's thesis is sharp: the triangle is a
coordinate system, not a tradeoff curve --- and \cref{sec:frontier-2026}'s
snapshot is the localisation, not the resolution.

\subsection{The six contradictions}
\label{sec:contradictions-six}

\paragraph{(C1) PRM utility post-R1.}
The PRM800K-era headline ``PRM > ORM at matched $N$''
\citep[\S 4]{lightman2023-prm800k}\footnote{KB excerpt:
\texttt{kb/excerpts/lightman2023-prm800k.md\#sec-4-2-aggregation}.} was
established on a non-reasoning-trained generator. Once the generator is
itself an RLVR-trained reasoner, the comparison shifts: the ORM in the
original setup is replaced by a zero-error symbolic verifier
(\cref{sec:r1-protocol}), and the PRM must now beat the policy's own
internal verification. \Citet{thinkprm2025} report inference-time gains
from generative PRMs over R1-distill outputs;
\citet{zheng2025-prm-survey} catalogue the 2025 picture as ``most PRMs
improve test-time-compute performance until the base model becomes a
strong reasoning model itself.'' R1 ships with no PRM in the headline
RL loop \citep[\S 2.2]{deepseek-r1}. The unresolved question splits
cleanly: PRMs survive as inference-time rerankers of best-of-$N$ and as
guidance signals for tree search (\cref{sec:tree-search,sec:prm-survival-inference});
their survival as an RL-loop \emph{training} signal beyond
verifier-only RLVR is contested
(\cref{sec:prm-survival-training}). Closure requires a controlled
PRM-shaped GRPO vs verifier-only GRPO comparison at matched
post-training compute on a reasoning-trained base; no published 2025--26
study has run it.

\paragraph{(C2) CoT faithfulness scaling.}
Two empirical waves are in tension. \Citet{arcuschin2025-cot-wild}
report a measurable in-the-wild rationalisation rate of $0.37\%$ on
naturalistic prompts for R1 \citep[abstract]{arcuschin2025-cot-wild},
plus an ``unfaithful illogical shortcut'' failure mode where trace
tokens are emitted in the right format but the answer is not in fact a
function of them (\cref{sec:faith-arcuschin}). \Citet{mehta2026-faithfulness-scaling}
measure non-thinking versus thinking configurations across frontier
models and find a per-rate gap of two orders of magnitude favouring
thinking models, but they also report that on at least one frontier
model accuracy drops while faithfulness rises under self-consistency
\citep[\S 4]{mehta2026-faithfulness-scaling}\footnote{KB note:
\texttt{kb/notes/reasoning/chain-of-thought.md\#4.1}.} --- a pattern
incompatible with the simple reading that traces are a causal substrate
at all scales. \Citet{shen2025-faithcot-bench} propose FaithCoT-Bench
as a dedicated faithfulness benchmark
(\cref{sec:faith-bench}). The unresolved question: does faithfulness
improve, plateau, or degrade with model scale and reasoning-training
intensity? The published 2026 ablations cannot separate long-CoT SFT,
RLVR post-training, and general scale at fixed others
(\cref{sec:faith-mehta}); the controlled ablation has not been run at
frontier scale, and Mehta's non-monotonic per-model results suggest
faithfulness sits on a separate curve from accuracy. Closure requires
a one-axis-at-a-time factorial at frontier scale, which the disclosure
regime of \cref{sec:frontier-2026} currently precludes for the closed
rows.

\paragraph{(C3) RLVR generalisation vs memorisation.}
The R1 narrative \citep[\S 2.3]{deepseek-r1}\footnote{KB excerpt:
\texttt{kb/excerpts/deepseek-r1.md\#sec-2-3-aime}.} reads the
$15.6\%\!\to\!77.9\%$ AIME 2024 pass@1 climb as RLVR-elicited
reasoning capability. \Citet{rlvr-limits-2025}%
\deepencite{rlvr-limits-2025 §3-§5, KB-excerpt-pending}
read it differently: the
group-mean baseline of GRPO (\cref{sec:rlvr-grpo-objective}) bounds
RLVR's reach by the base policy's pass@$K$ at \emph{any} $K$, not by
its pass@$1$, and so the climb is ``the base model could already solve
$77.9\%$ of these problems at some sample budget; GRPO transferred
that to pass@1.'' Under this reading RLVR \emph{sharpens} the policy
within base-model competency but does not \emph{expand} the solvable
set. The 2026 picture is contested rather than settled: the
\citeauthor{rlvr-limits-2025} reading is one camp's first-order
synthesis (the trace-growth result of \citet[\S 2.2]{deepseek-r1},
response length $\sim 200 \to \sim 10{,}000$ tokens, is not reconciled
with the capability-ceiling claim, and one cleanly-stated form of the
camp's synthesis is that ``RLVR teaches \emph{trace-length use} (when
to spend more tokens) without teaching new primitives''
\citep[\S 5]{rlvr-limits-2025}), but at least two camps have published
2025-Q3+ counter-evidence. \Citet{wen2025-cot-passk} introduce
CoT-Pass@K --- a metric that scores both the final answer and the
chain-of-thought derivation --- and report a consistent and
significant gap between RLVR-trained and base models on
AIME 2024/2025 across all $K$ up to $1024$, in contrast to standard
Pass@K where the gap closes; this directly challenges the
boundary-collapse reading on the metric most exposed to whether the
chain itself is novel. \Citet{rl-plus-2025} explicitly frame Yue's
findings as ``capability boundary collapse'' and propose
hybrid-policy optimization that --- on their reported runs ---
maintains an advantage over both the base model and GRPO as $K$ grows;
the field broadly agrees the boundary issue Yue diagnoses is real, but
disagrees about whether it is intrinsic to RLVR or tractable under
better optimisation. Whether either camp's synthesis exhausts the
picture or merely captures the dominant term on its preferred
benchmark is open. Closure requires
either a published pass@$K$-vs-base-model audit at large $K$ on R1's
training distribution, or a controlled experiment where RL is run on a
base whose pass@$K$ on the target benchmark is provably zero at
$K = 10^{6}$.

\paragraph{(C4) Search-vs-RL under fixed compute.}
\Citet[\S 6]{snell2024} establish a $14\times$ pretraining-vs-inference
substitution on the easy/intermediate MATH bin under PaLM~2-S\*.
\Citet[\S 2.3]{deepseek-r1} establish that GRPO on DeepSeek-V3-Base
converts ${\sim}10^{4}$ training steps into $+62.3$ AIME 2024 pass@1
points with self-consistency adding a further $+8.8$ at $N\!=\!64$.
\Citet[\S 4]{dapo2025} establish that DAPO halves $C_\text{RL}$ at
matched post-training accuracy. Each of the three results is an anchor
anecdote on a different base model, benchmark, and verifier; no
published study runs a three-channel FLOPs-matched experiment holding
$C_\text{total} = C_\text{pre} + C_\text{RL} + C_\text{test}$ fixed
and varying the allocation, and no parametric fit with the structural
status of the Chinchilla law $L(N, D) = E + A/N^\alpha + B/D^\beta$
exists for the joint $(C_\text{train}, C_\text{RL}, C_\text{test})$
frontier (\cref{sec:svr-identity}). Whether such a clean closed form
exists \emph{at all} is the open frontier question. The contingencies
of \cref{sec:svr-benchmark-specific} suggest the answer may be ``only
within a fixed prompt-difficulty distribution, verifier setting, base
model, and deployment regime,'' which is a much weaker statement than
the universal scaling law that would generalise Chinchilla to the
inference side. The operating-point heuristics of
\cref{sec:svr-heuristics} are what is locally available; a global
exchange rate is not.

\paragraph{(C5) AIME 2024 contamination.}
The headline R1 jump (\cref{sec:r1-contamination}) attributes the
$15.6\%\!\to\!77.9\%$ AIME 2024 climb to RLVR-elicited reasoning. The
competing reading: AIME-style problems and their solutions are widely
indexed on the open web --- problem-set archives, contest-prep
textbooks, math forums, AoPS-style community posts --- and \emph{some
fraction} of the pretraining-corpus distribution overlaps with the
test distribution. Under that reading, part of the headline gain is
recall of memorised solutions surfaced through long-CoT self-prompting,
not novel reasoning composed under RL. The arxiv preprint of
\citet{deepseek-r1} did not report a contamination audit on the AIME
2024 train/test boundary; standard contamination-detection
methodologies (n-gram overlap on solutions, near-duplicate matching,
decontamination-by-substring) all have known false-negative modes
against post-hoc-paraphrased problems and arithmetic isomorphisms, and
as of 2026-Q2 no \emph{third-party} study has reproduced an audit on
R1's pretraining mix (\cref{sec:r1-contamination}). The contradiction
interacts non-trivially with (C3): if the base model's pass@$K$ at
moderate $K$ is already high because of contamination, then the
\citeauthor{rlvr-limits-2025} reading becomes near-tautological on
this benchmark. Cross-paper: Paper~5~\S contamination is the
load-bearing treatment of the methodology debate itself --- which
methodologies disagree on what counts as contamination, and what
contamination-robust evaluation looks like.

\textbf{Update (post-2025-09).} The \emph{Nature} publication of R1
\citep{deepseek-r1-nature-2025} narrows but does not close this
contradiction. The peer-reviewed version adds (i) an n-gram-overlap
decontamination audit on the pre- and post-training data of R1, and
(ii) a control experiment that re-runs the GRPO recipe on Qwen2-7B ---
a base model released June 2024, before any advanced reasoner shipped
--- and reproduces similar capability emergence
\citep{deepseek-r1-nature-2025}. The Nature reviewer-response thread,
published alongside the paper, also addresses the sister
``distilled-from-OpenAI-outputs'' concern. Two open seams remain after
this update: a \emph{third-party} (non-DeepSeek) audit on R1's
pretraining mix, and a contamination-robust AIME-class benchmark
adopted field-wide. The closure recommended in
\cref{sec:contradictions-closure} for (C5) therefore narrows to those
two items rather than ``any audit at all,'' but the deflationary
reading of the AIME headline is no longer the only honest reading; the
Qwen2-7B control is the strongest single piece of evidence against it
in the published 2026-Q2 literature.

\paragraph{(C6) MCTS vs beam at matched budget.}
\Citet[\S 5--6]{snell2024} report BFS-V (the level-order beam-with-PRM
of \cref{sec:tree-beam}) competitive with or outright better than
alternative search strategies at matched compute on MATH with
PaLM~2-S\* as the base model
\citep[\S 5.2]{snell2024}\footnote{KB excerpt:
\texttt{kb/excerpts/snell2024.md\#sec-5-2-search}.}: the extra MCTS
machinery (UCT bookkeeping, rollout-derived $Q(s)$, the PUCT
exploration constant $c_\text{puct}$) does not necessarily pay for
itself. \Citet{rstar-math2025} foreground MCTS as the central
mechanism by which 7B SLMs match o1-preview, and the AlphaZero
structural parallel of \cref{sec:tree-rest-mcts} runs entirely through
the UCT recurrence (\cref{sec:tree-mcts}). Reconciling the two: Snell's
evaluation uses a non-reasoning-trained base and a single round of
search; rStar-Math co-trains MCTS with the policy and PRM across four
rounds of self-evolution \citep[abstract]{rstar-math2025}. The
``MCTS'' that wins in rStar-Math is a different object from the
``MCTS'' that ties in Snell. Whether the MCTS-vs-BFS gap reverses on
reasoning-trained base models, on harder benchmarks, or under
multi-round co-evolution is not closed in 2026; the
\citet{wei2025-tree-search-survey} taxonomy catalogues the
disagreement without resolving
it.\deepencite{wei2025-tree-search-survey \S5 (resolution discussion;
excerpt file pending)} Closure requires a search-strategy comparison
at matched budget on a reasoning-trained base under multi-round
co-evolution; published 2025--26 work runs at most two of those three
controls simultaneously.

\subsection{What closure would require}
\label{sec:contradictions-closure}

\begin{contradiction}
\textbf{The six contradictions share a methodological obstacle.}
Every closure listed above is gated by an experiment whose 2026-Q2
analogue does not exist in the public literature: a controlled,
matched-compute factorial on a frontier-scale reasoning-trained base,
with the closed-frontier vendors disclosing enough of their training
and inference pipeline to make a third-party replication possible.
The disclosure ceiling of \cref{sec:frontier-2026} is the binding
constraint: the published evidence base on R1 is the only frontier-row
in \cref{tab:frontier-2026} where the absence of a PRM in the
headline reward is positively attested rather than not-disclosed
\citep[\S 2.2]{deepseek-r1}. Two readings of the apparent
column-uniformity of frontier reasoning models are both compatible
with the 2026-Q2 evidence: \emph{(i)} the field has converged on
verifier-only RL plus long-CoT plus thinking-budget knobs as the
dominant pattern, or \emph{(ii)} the closed-frontier vendors are
running undisclosed PRM or search machinery and the disclosure ceiling
produces the apparent convergence as a measurement artefact. The
six contradictions are not just empirically open --- they are open
\emph{conditional on} a disclosure regime that the field cannot
unilaterally change.
\end{contradiction}

The closures tabulated by contradiction:
\emph{(C1)} a PRM-shaped GRPO vs verifier-only GRPO comparison at
matched post-training compute on a reasoning-trained base;
\emph{(C2)} a one-axis-at-a-time factorial separating long-CoT SFT,
RLVR post-training, and general scale at fixed others, run at frontier
scale and on a faithfulness-stable benchmark;
\emph{(C3)} a published pass@$K$-vs-base-model audit at large $K$ on
R1's training distribution, or RL on a base whose pass@$K$ at $K = 10^{6}$
is provably zero on the target benchmark;
\emph{(C4)} a three-channel FLOPs-matched experiment varying the
allocation of $C_\text{pre}$, $C_\text{RL}$, $C_\text{test}$ at fixed
$C_\text{total}$, with at least three base models and three benchmark
families to separate base-model from prompt-distribution effects;
\emph{(C5)} a third-party contamination audit on R1's pretraining mix
plus a contamination-robust AIME-class evaluation methodology
established cross-paper in Paper~5;
\emph{(C6)} a search-strategy comparison at matched budget on a
reasoning-trained base under multi-round co-evolution.
None of the six closures is a closed-form derivation; all are
empirical experiments whose run cost rivals frontier post-training.
The 2026-Q2 reasonable expectation is that the closures will arrive
piecemeal, on different benchmarks, with caveats that limit
generalisation --- not as a unified resolution.

\subsection{What is settled}
\label{sec:contradictions-settled}

\begin{intuition}
The contradictions are about \emph{margins}, not foundations. What
the 2025--26 literature does establish is structural: the three legs
of the triangle exist as separable mechanisms; each composes with the
others; the productionised exposure of inference-time compute
(thinking-budget knobs, reasoning-effort levels) has converged
(\cref{sec:frontier-2026}); and the empirical anchors on each leg
\citep{snell2024,lightman2023-prm800k,shao2024,deepseek-r1,dapo2025,s1-2025}
hold up across published replications. The triangle is real. The
contradictions concern which combination is optimal at which point, not
whether the legs themselves are real.
\end{intuition}

The settled-vs-contested split, leg by leg, parallels the
agreement-and-divergence summary of
\cref{sec:frontier-2026}. Settled: \emph{every} frontier-2026 reasoning
entry ships RLVR plus long-CoT, \emph{none} ship a process reward
model in the headline RL loop, and the productionised exposure of
inference-time compute has converged on a small number of API-shaped
``thinking budget'' surfaces. Contested: the contamination-vs-elicitation
attribution of the AIME-class headlines (C3, C5), the form of the
inference-time scaling law (the $L(\theta, N)$ functional form open
question of \cref{sec:snell-open}, partially overlapping with C4), the
faithfulness behaviour of long-CoT under scale (C2), and the
PRM-vs-verifier and beam-vs-MCTS implementation choices (C1, C6).

\subsection{Closing thesis}
\label{sec:contradictions-closing}

The reasoning gains observed between 2022 and 2026 are the joint
product of inference-time compute scaling, verifier signals, and RL on
verifiable rewards. The empirical evidence that these three legs are
co-designed --- not three independent gains shipped in convoy --- is
strongest on the open-weight rows of \cref{tab:frontier-2026},
weakest on the closed-frontier rows where the disclosure ceiling
\citep[\S 2.2]{deepseek-r1}\footnote{KB excerpt:
\texttt{kb/excerpts/deepseek-r1.md\#sec-2-2-reward}.} restricts what
can be positively attested. The six contradictions of
\cref{sec:contradictions-six} are not anomalies to be explained away;
they are the live frontier of the 2026 reasoning research programme.

The framing this paper has pushed throughout is that the productive
2026 question is not ``is reasoning real'' --- the triangle's
empirical effects on AIME, MATH, GPQA, and the agentic-benchmarks of
\cref{sec:frontier-2026} are not in dispute --- but ``which combination
of legs is optimal at which $(C_\text{train}, C_\text{RL}, C_\text{test})$
operating point, for which prompt-difficulty distribution, on which
base model, under which deployment latency budget?'' The 2025--26
literature has not closed any of those questions, and the closures
listed in \cref{sec:contradictions-closure} are empirical experiments
of frontier-scale cost gated by a disclosure regime the field does not
yet operate under. The reasonable forecast is that the answer will not
be closed-form for at least another model generation. What the rest of
the series can build on is the coordinate system this paper has named
--- inference compute, verifier signal, RL on verifiable reward --- and
the localisation of each open question on it. Paper~4
\S faithfulness and Paper~5 \S AIME-contamination both backreference
this section's (C2) and (C5) directly; Paper~2's RLHF-to-RLVR
transition treats (C3) as the hinge between learned and rule-based
reward sources.

The triangle is real and co-designed in frontier 2026 reasoning
models. The open questions are about margins and operating points,
not foundations. They will outlast the next model generation.
